\noindent
Let us introduce a notation: given $n \geq c \geq 0$
{\it an $(A, n, c)$-module}
is a finite $A$-module $M$ which is annihilated by $I^n$
and which as an $A/I^n$-module is $I^c/I^n$-projective, see
More on Algebra, Section \ref{more-algebra-section-near-projective}.

\medskip\noindent
We will use the following abuse of notation: given an $A$-module $M$
we denote $p^*M$ the quasi-coherent module gotten by pulling back by $p$
the quasi-coherent module $\widetilde{M}$ on $\Spec(A)$ associated to $M$.
For example we have $\mathcal{O}_{Y_n} = p^*(A/I^n)$.
For a short exact sequence $0 \to K \to L \to M \to 0$ of $A$-modules
we obtain an exact sequence
$$
p^*K \to p^*L \to p^*M \to 0
$$
as $\widetilde{\ }$ is an exact functor and $p^*$ is a right exact functor.
